\documentclass[../main.tex]{subfiles}

\begin{document}
\subsection{\textit{Learn Stack Layout}}
Saat jam-jam terakhir di perusahaan X, Holil dan Gilang ingin mencoba melakukan eksperimen yaitu melakukan \textit{debugging} pada struktur \textit{address} dari fungsi yang diinput di sebuah bahasa C yang di {\compile} dengan 64-bit. Holil dan Gilang melakukan \textit{debugging} untuk melihat sebuah fungsi yang memiliki atau tidak memiliki \textit{argument}. Ternyata setelah holil dan gilang belajar tentang \textit{register}, mereka mengetahui bahwa pada 64-bit, \textit{positional argument} diurutkan mulai dengan RDI, RSI, RSI, RDX, RCX, R8, R9. Holil mendapat penjelasan berikut:

Ketika fungsi memiliki paramater, maka fungsi tersebut dipanggil dengan \textit{argument} dan jumlah argument yang masuk akan diurutkan sesuai dengan ketentuan \textit{register} pada 64-bit.

\inputfigure{res/m6i1.png}{0.5}{Visualisasi Register Pada Fungsi 2 Argumen}{fig:m6:i1}
\vspace*{1.6ex}
\inputfigure{res/m6i2.png}{0.5}{Visualisasi Register Pada Fungsi 3 Argumen}{fig:m6:i2}
\vspace{1ex}
\subsection{\textit{Dungeon Crawler} RPG}

Setelah pulang dari perusahaan X, Holil sampai di rumah dan ingin memainkan game \textit{roguelike online} yang pernah dia mainkan. \textit{Roguelike} adalah \textit{subgenre} permainan video permainan peran yang ditandai dengan \textit{dungeon crawl} melalui level yang dihasilkan secara prosedural, \textit{gameplay} berbasis giliran, grafis berbasis ubin dan karakter pemain dengan kematian permanen. Karena merasa akan bosan dan hambar jika dimainkan sendirian, Holil pun mengajak temannya, Gilang, seorang pengembang \textit{game}, untuk bermain bersama. Setelah beberapa saat ingin memainkan game tersebut, Holil mendapat berita dari internet bahwa \textit{game online} tersebut sudah menutup server saat dia melamar kerja di perusahaan X. Holil pun sangat kecewa dengan penutupan server \textit{game} tersebut. Gilang, yang kebetulan seorang pengembang \textit{game}, berinisiatif mengajak Holil untuk membuat ulang \textit{game} tersebut agar bisa dimainkan. Dengan menggabungkan kemampuan pemrograman dari Holil dan kemampuan pengembangan game dari Gilang, mereka pun membuat ulang game tersebut menggunakan bahasa C++ yang dijalankan melalui \textit{command prompt}.

Pada bagian awal, \textit{game} akan menampilkan tampilan berikut:

\inputfigure{res/m6i18.png}{0.11}{Tampilan Penjelajahan Program \textit{Dungeon Crawl} RPG}{fig:m6:i18}

Untuk menjalankannya, \textit{player} dapat menggunakan perintah \citecode{next} atau \citecode{go} untuk menuju ke room selanjutnya dan perintah \textit{back} untuk kembali menuju ke \textit{room} sebelumnya.

\inputfigure{res/m6i19.png}{0.13}{Tampilan Bertarung Program \textit{Dungeon Crawl} RPG}{fig:m6:i19}

Jika terdapat lebih dari satu pilihan \textit{room}, \textit{game} akan memberikan pilihan untuk dipilih oleh \textit{player}. Jika \textit{room} selanjutnya terdapat musuh atau monster maka game akan langsung menuju ke mode bertarung.

\inputfigure{res/m6i20.png}{0.2}{Tampilan \textit{Item} Ditemukan Program \textit{Dungeon Crawl} RPG}{fig:m6:i20}

Jika terdapat suatu \textit{item} pada suatu \textit{room}, maka \textit{game} akan meminta \textit{player} memasukan nama \textit{player} yang akan mengambil \textit{item} tersebut, selanjutnya jika suatu \textit{item} diambil dan monster dikalahkan maka akan menampilkan pesan tidak ada sesuatu pada \textit{room} tersebut.

\inputfigure{res/m6i21.png}{0.21}{\textit{Boss} Terakhir Program \textit{Dungeon Crawl} RPG}{fig:m6:i21}

Objektif dari \textit{game} ini yaitu mengalahkan \textit{boss} terakhir bernama “Leviathan” pada \textit{room} “Abyssia Ocean”. Jika salah satu dari \textit{player} mati maka \textit{game} akan \textit{restart} dari awal sampai \textit{boss} terakhir dikalahkan. Program selesai Ketika berhasil mengalahkan \textit{boss} terakhir.

\end{document}
